\documentclass{report}
\usepackage{amsmath,amssymb}
\setlength{\parindent}{0mm}
\setlength{\parskip}{1em}
\begin{document}
\begin{center}
\rule{6in}{1pt} \
{\large Frederic Vi \\
{\bf Parallel Multigrid Preconditionner based on Automatic 3D Tetraedric Meshes}}

MINES ParisTech \\ CEMEF \\ 1 rue Claude Daunesse \\ CS 10207 \\ 06904 Sophia Antipolis Cedex \\ France
\\
{\tt frederic.vi@mines-paristech.fr}\\
Katia Mocellin\\
Hugues Digonnet\\
Lionel Fourment\end{center}

Multigrid methods are efficient for solving large sparse linear systems.
Geometric (GMG) and Algebraic Multigrid (AMG) have both their own
benefits and limitations. Combining the simplicity of AMG with the
efficiency of GMG lead us to the development of an Hybrid Multigrid
preconditionner. From an initial fine mesh we first build coarser
unnested meshes and then deduce interpolation and restriction matrices
based on interpolation of a mesh into a coarser one. We finally use
Galerkin relation on each level to create our coarser problem matrix.

We focus on the resolution of linear systems coming from solid mechanics
problems. More specifically we solve Stokes problem using a 3D tetraedric
mixed finite element formulation (P1+/P1) leading to a specific kind of
problem with the associated matrix [1]. To deal with complex geometries
and various material forming processes, we use unstructured mesh and need
the resolution method to be compatible with automatic remeshing. Coarser
meshes, interpolation and restriction operators and coarser matrices have
then to be automatically built and rebuilt.

As nowadays numerical simulations can not ignore the evolution of
computer architecture, our developments take parallel computing into
account. We manage to create well-balanced partitions and generate
coarser meshes with a similar partitioning, so that the communication
costs are minimized. The matrices (interpolation, restriction, coarser
problems) are computed in parallel and the resolution uses smoothers and
coarse grid solver suitable for parallel computations.


[1] T. Coupez, S. Marie, From a direct solver to a parallel iterative
solver in 3-D forming simulation, International Journal of High
Performance Computing Applications (1997) 11: 277-285.


\end{document}
